\begin{definition}[The program of one process]
  \label{def:aba-programs}
  \lean{PLTS.ABA.GBCA.ByABDY.RoundRecord}\lean{PLTS.ABA.RoundLoopRecord}\lean{PLTS.ABA.ABDY.RoundRecordMap}\lean{PLTS.ABA.ABDY.ProcessRecord}\lean{PLTS.ABA.Composition.NetworkEvent}\lean{PLTS.ABA.Composition.ExtendedLabel}\lean{PLTS.ABA.Implementation.networkEventLabels}\lean{PLTS.ABA.ABDY.RoundStep}\lean{PLTS.ABA.ABDY.ABAProgramStep}\lean{PLTS.ABA.ABDY.ABAProgram}\leanok
  \uses{def:aba-labels, def:aba-round-messages, def:aba-round-loop-state, def:aba-implementation, def:relabel, def:synchronised-product}
  Process $j$ runs Definition~\ref{def:aba-implementation} at ABDY22's implementation of Definition~\ref{def:aba-round-messages}: the round message type is $\mathrm{GBCA.ByABDY.Message}$, the round record is $\mathrm{GBCA.ByABDY.RoundRecord}(n)$, and the round's transitions are the fourteen of \leandecl{PLTS.ABA.ABDY.RoundStep} --- the eight multicasts of those levels, the delivery that files a received message, the call and its loop, and the three returns.
  The twenty-nine transitions around them are the shared ones.
  The state of process $j$ is
  \[
    \mathrm{ProcessRecord}(n) \;=\;
      \mathrm{RoundLoopRecord}(n) \times \mathrm{RoundRecordMap}(n),
  \]
  its round-loop record --- the control record of Definition~\ref{def:aba-round-loop-state} together with the DECIDED received sets $decidedDelivered(k)$ --- beside its round's record.
  The round's record holds a finite map $roundRecords$ from rounds to round records of Definition~\ref{def:aba-round-messages}, each with its received sets $received(k)$, and a flag $terminated$ (D22).
  The round record of round $r$ is read $\mathrm{roundRecord}\,r$: the retained record where $j$ has touched round $r$, the initial round record elsewhere.
  The map is finite by its type, so $j$ holds finitely many round records at every point of a run, and the round advance resets none of them.
  Every round's transition is accordingly round-indexed rather than round-guarded: it reads and writes the round record of the round its own label tags, whichever round the round loop is in, under the guards of that round record and $terminated = \mathrm{false}$.
  Process $j$ therefore answers the messages of a round its round loop has passed, and a round delivery of any round is filed.
  The transition $terminate$ fires on $j$'s own return and $2f+1$ DECIDED receipts, and its effect is to set $terminated$ and nothing else, so the retained round records remain unchanged from that point.
  The DECIDED transitions carry no termination guard, so a process that has terminated keeps relaying the payloads it holds.
  The return $retABA$ and the DECIDED relay are guarded by the round-loop record holding an input, so a process that was never called neither returns nor relays (D8).
  The transition $gbcaCallLoop$ reads the round record of round $r$ without writing it, requiring that the record holds an input, and carries no termination guard.
  A process that has terminated at phase $toCallG$ over an uncalled round record therefore has a transition on neither call label and takes no further round-loop step, an excluded region this definition accepts.
  There is no outbox and no copy of the corrupted set or its budget, so every guard of the automaton $\mathrm{ABAProgram}(P,j)$ reads $j$'s own record.
  A corruption replaces $j$'s program (D23).
  The round-loop record carries the flag \leandecl{PLTS.ABA.RoundLoopRecord.corrupted}, written by $j$'s own half of the corruption broadcast, \leandeclnamed{failSelf}{PLTS.ABA.Implementation.ProgramStep.failSelf}; every other transition above is enabled only while that flag is down, so the record is unchanged from the corruption and $terminate$ in particular never fires afterwards.
  In place of those transitions the replaced program is the single self-loop \leandeclnamed{corruptedIdle}{PLTS.ABA.Implementation.ProgramStep.corruptedIdle}, taken on every label other than $\tau$ and the labels of $\mathrm{actsAt}\,j$ --- the labels on which $j$ would act on its own sub-protocol messages --- where it has no transition at all.
  Those messages enter through the handshake transitions of D11 and the injections $\mathrm{byzantineGBCA}$ and $\mathrm{byzantineDecided}$.
  The visible return is the exception: the network's own $\mathrm{retByzantine}$ transition requires $j \in F$ and no DECIDED evidence, and the self-loop is its other half.
  It runs over the extended alphabet $\mathrm{ExtendedLabel}(n, M) = \mathrm{Label}(n) \oplus \mathrm{NetworkEvent}(n, M)$, over the type $M$ of the messages a graded-agreement round exchanges, where $\mathrm{NetworkEvent}$ names the rendezvous the shared alphabet cannot: the round multicast $\mathrm{gbcaSend}(r,j,m)$ and the round delivery $\mathrm{gbcaDeliver}(r,i,j,m)$; the DECIDED relay $\mathrm{decidedSend}(j,b)$ and the DECIDED delivery $\mathrm{decidedDeliver}(i,j,b)$; the fused coin return $\mathrm{retWPublish}(r,id,c,b)$, which carries beside the coin the payload to be published (D10); the graded-agreement call $\mathrm{gbcaCallLoop}(r,id,b)$ against an already-opened round record; and the five Byzantine handshake transitions $\mathrm{byzantineCallG}$, $\mathrm{byzantineCallGLoop}$, $\mathrm{byzantineRetG}$, $\mathrm{byzantineCallW}$ and $\mathrm{byzantineRetW}$ (D11).
  The set $\mathrm{networkEventLabels}(n)$ of these labels is what the composition hides, the silent label of $\mathrm{ExtendedLabel}(n, M)$ being that of $\mathrm{Label}(n)$.
  Only the round multicast and the round delivery carry a message, so every component written over the alphabet is written for every $M$; the ABDY chain takes $\mathrm{GBCA.ByABDY.Message}$ for it and the gather-based chain of Section~\ref{sec:aba-gather} takes the empty type.
  Every transition is Dirac, and $terminate$ is the one transition of the list that fires on $\tau$.
  While $j$'s program stands unreplaced, every label of either alphabet has a transition of process $j$, the participant's or an idle self-loop, with two exceptions: $\mathrm{byzantineCallG}(r,k,b)$ and $\mathrm{byzantineRetG}(r,k,out,bnd)$ have no transition at the process $k$ they name (D11, D22), and a corrupted process's round messages enter through the network's own $\mathrm{byzantineGBCA}$ injection instead.
  Once the program is replaced the labels of $\mathrm{actsAt}\,j$ join those two (D23).
  Full synchronisation therefore lets a label's owner move while every other program stands.
  The full list of transitions, forty-three of them --- twenty-nine shared, fourteen of ABDY22's: \leandecl{PLTS.ABA.ABDY.ABAProgram}.
\end{definition}
